% Part: second-order-logic
% Chapter: sol-and-sets
% Section: power-of-continuum

\documentclass[../../../include/open-logic-section]{subfiles}

\begin{document}

\olfileid{sol}{set}{pow}

\olsection{সতত সমষ্টির অঙ্কবাচকতা}

\begin{explain}
দ্বিতীয়-ক্রমের যুক্তিবিদ্যায় আমরা !!{domain}-এর উপসেটের উপর পরিমাণায়ন
করতে পারি, কিন্তু !!{domain}-এর উপসেটের সেটের উপর পারি না। তা সরাসরি করতে
হলে \emph{তৃতীয়-ক্রমের} যুক্তিবিদ্যা দরকার হতো। যেমন, কোনো সেটের ঘাত সেট
থেকে সেটটিতেই কোনো !!{injective} অপেক্ষক নেই---কান্টরের উপপাদ্যের এই
বক্তব্যটি বলতে চাইলে আমরা হয়তো এভাবে সূত্রবদ্ধ করার চেষ্টা করতাম: ``প্রতিটি
সেট~$X$ ও প্রতিটি সেট~$P$-এর জন্য, $P$ যদি~$X$-এর ঘাত সেট হয়, তবে
$\cardle{P}{X}$ নয়।'' আর~$P$ যে~$X$-এর ঘাত সেট, তা বলতে আনুষ্ঠানিকভাবে
বলতে হতো যে~$P$-এর !!{element}s ঠিক~$X$-এর উপসেটগুলিই; তাই
$\lforall[Y][(P(Y) \liff Y \subseteq X)]$-এর মতো কিছু লাগত। সমস্যাটি
$P(Y)$-তে: এটি দ্বিতীয়-ক্রমের যুক্তিবিদ্যার কোনো !!a{formula} নয়, কারণ
$P$-এর মতো এক-স্থানীয় সম্পর্ক-চলরাশির আর্গুমেন্ট কেবল পদই হতে পারে।

তবে সংজ্ঞাক্ষেত্র যথেষ্ট বড় হলে সেটের সেটের উপর পরিমাণায়নকে আমরা
\emph{অনুকরণ} করতে পারি। ধারণাটি হলো, দ্বি-স্থানীয় সম্পর্ক~$R$ যে
!!{domain}-এর !!{element}s-কে !!{domain}-এর !!{element}s-এর সঙ্গে
সম্পর্কিত করে, সেই তথ্য ব্যবহার করা। এমন একটি~$R$ দেওয়া থাকলে, কোনো~$x$
যে সব !!{element}s-এর সঙ্গে $R$-সম্পর্কিত, তাদের একত্র করতে পারি:
$\Setabs{y \in \Domain{M}}{R(x,y)}$ হলো~$x$-এর দ্বারা ``সংকেতায়িত''
সেট। বিপরীত দিকে, $Z \subseteq \Pow{\Domain{M}}$ যদি
~$\Domain{M}$-এর উপসেটের কোনো সমষ্টি হয়, এবং~$Z$-তে যত সেট আছে
~$\Domain{M}$-এ অন্তত ততগুলি !!{element}s থাকে, তাহলে এমন একটি সম্পর্ক
$R \subseteq \Domain{M}^2$-ও আছে যার সাহায্যে প্রতিটি $Y \in Z$-কে
কোনো~$x$,~$R$-এর সাহায্যে সংকেতায়িত করে।
\end{explain}

\begin{defn}
$R \subseteq \Domain{M}^2$ হলে, \emph{$x$, $R$-এর সাহায্যে}
$\Setabs{y \in \Domain{M}}{R(x, y)}$-কে \emph{সংকেতায়িত করে}।
\end{defn}

কোনো !!a{element}~$x \in \Domain{M}$ যদি $R$-এর সাহায্যে কোনো সেট $Z
\subseteq \Domain{M}$-কে সংকেতায়িত করে, তাহলে একটি সেট $Y \subseteq
\Domain{M}$ সেটের একটি সমষ্টিকে সংকেতায়িত করে, যথা~$Y$-এর !!{element}s
যে সেটগুলিকে সংকেতায়িত করে। তাই কোনো সেট~$Y$, $R$-এর সাহায্যে
$\Pow{X}$-কে সংকেতায়িত করতে পারে। এমনটি হয় যদি এবং কেবল যদি প্রতিটি $Z
\subseteq X$-এর জন্য কোনো $x \in Y$,~$R$-এর সাহায্যে~$Z$-কে সংকেতায়িত
করে, এবং প্রতিটি $x \in Y$, কোনো~$Z \subseteq X$-কে~$R$-এর সাহায্যে
সংকেতায়িত করে।

\begin{prop}
!!{formula}
  \[
  \fn{Codes}(x, R, Z) \ident \lforall[y][(Z(y) \liff
    R(x, y))]
  \]
প্রকাশ করে যে $s(x)$,~$s(R)$-এর সাহায্যে~$s(Z)$-কে সংকেতায়িত করে।
নিচের !!{formula}টি
\begin{multline*}
  \fn{Pow}(Y, R, X) \ident \\
  \lforall[Z][(Z \subseteq X \lif \lexists[x][(Y(x) \land
    \fn{Codes}(x, R, Z))])] \land {} \\ \lforall[x][(Y(x) \lif
  \lforall[Z][(\fn{Codes}(x, R, Z) \lif Z \subseteq X)]]
\end{multline*}
প্রকাশ করে যে $s(Y)$,~$s(R)$-এর সাহায্যে~$s(X)$-এর ঘাত সেটকে সংকেতায়িত
করে; অর্থাৎ, $s(Y)$-এর উপাদানগুলি $s(R)$-এর সাহায্যে ঠিক~$s(X)$-এর
উপসেটগুলিকেই সংকেতায়িত করে।
\end{prop}

\begin{explain}
এই কৌশলের সাহায্যে উপসেটগুলির উপর পরিমাণায়নের বদলে তাদের সংকেতগুলির উপর
পরিমাণায়ন করে ঘাত সেট সম্বন্ধীয় বক্তব্য প্রকাশ করা যায়। যেমন, এবার
কান্টরের উপপাদ্য এভাবে প্রকাশ করা যায়:~$X$-এর ঘাত সেটকে সংকেতায়িত করে
এমন কোনো সম্পর্কের সংকেতগুলির সংজ্ঞাক্ষেত্র থেকে~$X$-তেই কোনো
!!{injective} অপেক্ষক নেই।
\end{explain}

\begin{prop}
নিচের বাক্যটি
\begin{multline*}
  \lforall[X][\lforall[Y][\lforall[R][(\fn{Pow}(Y, R, X) \lif \\
      \lnot \lexists[u][(\lforall[x][\lforall[y][((Y(x) \land Y(y) \land
            \eq[u(x)][u(y)]) \lif \eq[x][y])]] \land {}\\
        \lforall[x][(Y(x) \lif X(u(x)))])])]]]
\end{multline*}
বৈধ।
% BN-SRC-266 TARGET-CORRECTION-BEGIN
\par\smallskip\noindent\textnormal{\small\textbf{উৎস-সংশোধন (BN-SRC-266):} হিমায়িত সূত্রে~$Y$ থেকে~$X$-তে একৈক অপেক্ষকের অস্তিত্ব অস্বীকার করার বদলে~$u$-কে গোটা সংজ্ঞাক্ষেত্রে একৈক হতে বলা হয়েছে। বাংলা সূত্রে একৈকতার শর্ত~$Y$-এর উপাদানযুগলে সীমাবদ্ধ করা হয়েছে, যাতে বাক্যটি ঠিক সংকেতসমষ্টি থেকে~$X$-তে একৈক অপেক্ষকের অনস্তিত্ব প্রকাশ করে।} % BN-SRC-266 TARGET-CORRECTION-END
\end{prop}

\begin{explain}
কোনো !!a{denumerable} সেটের ঘাত সেট !!{nonenumerable}; তাই তার অঙ্কবাচকতা
যেকোনো !!{denumerable} সেটের অঙ্কবাচকতার (অর্থাৎ~$\aleph_0$-র) চেয়ে
বড়। $\Pow{\Nat}$-এর আকারকে ``সতত সমষ্টির অঙ্কবাচকতা'' বলা হয়, কারণ
এটি বাস্তব সংখ্যারেখার বিন্দুগুলির সেট~$\Real$-এর আকারের সমান।
!!{domain} যদি কোনো !!a{denumerable} সেটের ঘাত সেটকে সংকেতায়িত করার মতো
যথেষ্ট বড় হয়, তাহলে কোনো সেট সতত সমষ্টির আকারের---এ কথা এভাবে প্রকাশ
করা যায় যে, সেটটি~$\aleph_0$ আকারের কোনো সেট~$X$-এর ঘাত সেটকে
সংকেতায়িতকারী যেকোনো সেট~$Y$-এর সঙ্গে সমসংখ্যক। (সংজ্ঞাক্ষেত্র যথেষ্ট বড়
না হলে, অর্থাৎ তাতে~$\Real$-এর সঙ্গে সমসংখ্যক কোনো উপসেট না থাকলে,
$\Pow{X}$-কে সংকেতায়িত করে এমন কোনো সম্পর্কও থাকতে পারে না।)
\end{explain}

\begin{prop}
$\cardle{\Real}{\Domain{M}}$ হলে নিচের !!{formula}টি
\begin{multline*}
\fn{Cont}(Y) \ident
\lexists[X][\lexists[R][((\fn{Aleph}_0(X) \land
      \fn{Pow}(Y, R, X)) \land \\
      \lforall[x][\lforall[y][((Y(x) \land {}
      Y(y) \land \lforall[z][R(x,z) \liff R(y,z)]) \lif \eq[x][y])]])]]
\end{multline*}
প্রকাশ করে যে $\cardeq{s(Y)}{\Real}$।
\end{prop}

\begin{proof}
  $\fn{Pow}(Y, R, X)$ প্রকাশ করে যে $s(Y)$,~$s(R)$-এর সাহায্যে
  $s(X)$-এর ঘাত সেটকে সংকেতায়িত করে; আর $\fn{Aleph}_0(X)$ বলে যে
  সেটটি অসীম তালিকায়নযোগ্য। তাই $s(Y)$ অন্তত সতত সমষ্টির অঙ্কবাচকতার
  সমান বড়, যদিও সেটি আরও বড় হতে পারে ($s(Y)$-এর একাধিক উপাদান যদি
  $X$-এর একই উপসেটকে সংকেতায়িত করে)। শেষ সংযোজ্যটি সেই সম্ভাবনা
  বাদ দেয়: সেটি চায়, $s(R)$-এর মাধ্যমে~$s(Y)$-এর উপাদান ও~$s(X)$-এর
  উপসেটগুলির মধ্যেকার সম্বন্ধটি !!{injective} হোক।
% BN-SRC-264 TARGET-CORRECTION-BEGIN
\par\smallskip\noindent\textnormal{\small\textbf{উৎস-সংশোধন (BN-SRC-264):} হিমায়িত প্রমাণের শেষ বাক্যে হঠাৎ~$s(Z)$-এর উপসেট বলা হয়েছে, যদিও এই প্রতিজ্ঞায়~$Z$ মুক্ত বা নির্দিষ্ট নয় এবং গোটা নির্মাণটি~$s(X)$-এর ঘাত সেট নিয়ে। বাংলা পাঠে~$s(X)$ পুনরুদ্ধার করা হয়েছে।} % BN-SRC-264 TARGET-CORRECTION-END
\end{proof}

\begin{prop}
$\cardeq{\Domain{M}}{\Real}$ যদি এবং কেবল যদি
\begin{multline*}
  \Sat{M}{\lexists[Y][(\fn{Cont}(Y) \land {}\\
    \lexists[u][(\lforall[x][Y(u(x))] \land
      \lforall[x][\lforall[y][(\eq[u(x)][u(y)] \lif \eq[x][y])]] \land {}\\
      \lforall[y][(Y(y) \lif \lexists[x][\eq[y][u(x)]])])])]}.
\end{multline*}
% BN-SRC-265 TARGET-CORRECTION-BEGIN
\par\smallskip\noindent\textnormal{\small\textbf{উৎস-সংশোধন (BN-SRC-265):} হিমায়িত সূত্রে~$Y$ কেবল ঘাত সেটের সব উপসেটকে সংকেতায়িত করে, কিন্তু একই উপসেটের বহু সংকেত বাদ দেয় না; আবার~$u$-র বিস্তৃতিতে~$Y$ আছে বলা হলেও~$u$-র সব মান~$Y$-তে বলা হয় না। ফলে সূত্রটি সংজ্ঞাক্ষেত্র ও সতত সমষ্টির মধ্যে দ্বৈকতা প্রকাশ করে না। বাংলা সূত্রে আগের সংজ্ঞা~$\fn{Cont}(Y)$ ব্যবহার করে~$Y$-এর ঠিক সতত-আকার নিশ্চিত করা হয়েছে এবং~$u$-কে সংজ্ঞাক্ষেত্র থেকে~$Y$-তে একৈক ও সমাপতিত করা হয়েছে।} % BN-SRC-265 TARGET-CORRECTION-END
\end{prop}

\begin{explain}
সতত সমষ্টি অনুমান হলো এই বক্তব্য যে, সতত সমষ্টির আকারই প্রথম
!!{nonenumerable} অঙ্কবাচকতা; অর্থাৎ~$\Pow{\Nat}$-এর আকার~$\aleph_1$।
\end{explain}

\begin{prop}
সতত সমষ্টি অনুমান সত্য যদি এবং কেবল যদি নিচের বাক্যটি
\[\fn{CH} \ident
\lforall[X][(\fn{Aleph}_1(X) \liff \fn{Cont}(X))]\] বৈধ।
\end{prop}

লক্ষ করো, সতত সমষ্টি অনুমান মিথ্যা যদি এবং কেবল যদি $\lnot \fn{CH}$
বৈধ---এ কথা সত্য নয়। কোনো !!a{enumerable} সংজ্ঞাক্ষেত্রে~$\aleph_1$
আকারের কোনো উপসেট নেই, আবার সতত সমষ্টির আকারেরও কোনো উপসেট নেই; তাই
কোনো !!a{enumerable} সংজ্ঞাক্ষেত্রে~$\fn{CH}$ সবসময় সত্য। তবে সতত সমষ্টি
অনুমান মিথ্যা হলে এবং কেবল তখনই বৈধ, এমন একটি আলাদা বাক্য দেওয়া যায়:

\begin{prop}
সতত সমষ্টি অনুমান মিথ্যা যদি এবং কেবল যদি নিচের বাক্যটি
\[\fn{NCH} \ident
\lforall[X][(\fn{Cont}(X) \lif \lexists[Y][(Y \subseteq X \land \lnot
    \fn{Count}(Y) \land \lnot \cardeq{X}{Y})])]\] বৈধ।
\end{prop}

\end{document}
